\documentclass[10pt,a4paper]{amsart}
\usepackage{amsmath,amssymb,amsthm,graphicx}
\usepackage{verbatim}
\newenvironment{reference}{\comment}{\endcomment}
%\newenvironment{reference}{}{}
\newenvironment{slogan}{\comment}{\endcomment}
\newenvironment{history}{\comment}{\endcomment}

% For commutative diagrams we use Xy-pic
\usepackage[all]{xy}

% We use 2cell for 2-commutative diagrams.
\xyoption{2cell}
% Theorem environments.
%
\theoremstyle{plain}
\newtheorem{theorem}[subsection]{Theorem}
\newtheorem{proposition}[subsection]{Proposition}
\newtheorem{lemma}[subsection]{Lemma}

\theoremstyle{definition}
\newtheorem{definition}[subsection]{Definition}
\newtheorem{example}[subsection]{Example}
\newtheorem{exercise}[subsection]{Exercise}
\newtheorem{situation}[subsection]{Situation}

\theoremstyle{remark}
\newtheorem{remark}[subsection]{Remark}
\newtheorem{remarks}[subsection]{Remarks}

\numberwithin{equation}{subsection}

% Macros
%
\def\lim{\mathop{\mathrm{lim}}\nolimits}
\def\colim{\mathop{\mathrm{colim}}\nolimits}
\def\Spec{\mathop{\mathrm{Spec}}}
\def\Hom{\mathop{\mathrm{Hom}}\nolimits}
\def\Ext{\mathop{\mathrm{Ext}}\nolimits}
\def\SheafHom{\mathop{\mathcal{H}\!\mathit{om}}\nolimits}
\def\SheafExt{\mathop{\mathcal{E}\!\mathit{xt}}\nolimits}
\def\Sch{\mathit{Sch}}
\def\Mor{\mathop{\mathrm{Mor}}\nolimits}
\def\Ob{\mathop{\mathrm{Ob}}\nolimits}
\def\Sh{\mathop{\mathit{Sh}}\nolimits}
\def\NL{\mathop{N\!L}\nolimits}
\def\CH{\mathop{\mathrm{CH}}\nolimits}
\def\proetale{{pro\text{-}\acute{e}tale}}
\def\etale{{\acute{e}tale}}
\def\QCoh{\mathit{QCoh}}
\def\Ker{\mathop{\mathrm{Ker}}}
\def\Im{\mathop{\mathrm{Im}}}
\def\Coker{\mathop{\mathrm{Coker}}}
\def\Coim{\mathop{\mathrm{Coim}}}

% Boxtimes
%
\DeclareMathSymbol{\boxtimes}{\mathbin}{AMSa}{"02}

%
% Macros for moduli stacks/spaces
%
\def\QCohstack{\mathcal{QC}\!\mathit{oh}}
\def\Cohstack{\mathcal{C}\!\mathit{oh}}
\def\Spacesstack{\mathcal{S}\!\mathit{paces}}
\def\Quotfunctor{\mathrm{Quot}}
\def\Hilbfunctor{\mathrm{Hilb}}
\def\Curvesstack{\mathcal{C}\!\mathit{urves}}
\def\Polarizedstack{\mathcal{P}\!\mathit{olarized}}
\def\Complexesstack{\mathcal{C}\!\mathit{omplexes}}
% \Pic is the operator that assigns to X its picard group, usage \Pic(X)
% \Picardstack_{X/B} denotes the Picard stack of X over B
% \Picardfunctor_{X/B} denotes the Picard functor of X over B
\def\Pic{\mathop{\mathrm{Pic}}\nolimits}
\def\Picardstack{\mathcal{P}\!\mathit{ic}}
\def\Picardfunctor{\mathrm{Pic}}
\def\Deformationcategory{\mathcal{D}\!\mathit{ef}}
\usepackage{cleveref}
\crefname{originalequation}{original equation}{original equations}
\Crefname{originalequation}{Original equation}{Original equations}
\makeatletter
\newlabel{equation-restrict-modules}{{03LW}{source}{}{sourcepointer-1}{}}
\newlabel{equation-restrict-modules@cref}{{[originalequation][0][]03LW}{[1][0][]source}{}{}{}}
\makeatother
\crefname{originalreference}{original reference}{original references}
\Crefname{originalreference}{Original reference}{Original references}
\makeatletter
\newlabel{lemma-pullback-quasi-coherent}{{01BG}{source}{}{sourcepointer-2}{}}
\newlabel{lemma-pullback-quasi-coherent@cref}{{[originalreference][0][]01BG}{[1][0][]source}{}{}{}}
\makeatother
\crefname{originaldefinition}{original definition}{original definitions}
\Crefname{originaldefinition}{Original definition}{Original definitions}
\makeatletter
\newlabel{sites-modules-definition-global}{{03DE}{source}{}{sourcepointer-3}{}}
\newlabel{sites-modules-definition-global@cref}{{[originaldefinition][0][]03DE}{[1][0][]source}{}{}{}}
\makeatother
\crefname{originaldefinition}{original definition}{original definitions}
\Crefname{originaldefinition}{Original definition}{Original definitions}
\makeatletter
\newlabel{sites-modules-definition-site-local}{{03DL}{source}{}{sourcepointer-4}{}}
\newlabel{sites-modules-definition-site-local@cref}{{[originaldefinition][0][]03DL}{[1][0][]source}{}{}{}}
\makeatother
\crefname{originallemma}{original lemma}{original lemmas}
\Crefname{originallemma}{Original lemma}{Original lemmas}
\makeatletter
\newlabel{sites-modules-lemma-local-final-object}{{03DN}{source}{}{sourcepointer-5}{}}
\newlabel{sites-modules-lemma-local-final-object@cref}{{[originallemma][0][]03DN}{[1][0][]source}{}{}{}}
\makeatother
\title{Detection of quasi-coherent modules by pullback along an etale cover}
\author{The Stacks Project Authors}
\date{stacks-project:spaces-properties.tex:local-snapshot-20261012; GNU FDL 1.2 or later}
\begin{document}
\maketitle
\noindent\textit{Editorial note.} All delivered mathematical text is the author's own, copied from the pinned Stacks Project chapter file spaces-properties.tex; no statement or proof step was rewritten, reordered or removed.\par\medskip
\noindent\textit{Editorial note.} The counted proof is the author's complete argument: the trivial first implication, the descent of the hypothesis to a scheme chart by composing etale surjections, the production of global presentations on every object of the site from local presentations of the pullback, and the reduction of the last condition to affine charts.\par\medskip
\noindent\textit{Editorial note.} The label lemma-characterize-quasi-coherent names the counted statement, and the proof environment immediately following it in the source is delivered verbatim as its proof.\par\medskip
\noindent\textit{Editorial note.} The counted statement is carried by the author's own lemma environment in the chapter source.\par\medskip
\noindent\textit{Editorial note.} Pre-open gate: python operations/agent-authoring/tools/precard\_gate.py --source spaces-properties.tex --label lemma-characterize-quasi-coherent --start 5608 --end 5655 --local-id 28052 returned RESULT=OK (L28052 free, occupancy FREE, public interval clear). Candidate taken from sources/stacks-spaces-properties/candidates-halfA.json.\par\medskip
\noindent\textit{Editorial note.} The statement and proof are copied from the pinned chapter file; the cited inputs are kept as external pointers with their canonical Stacks tags.\par\medskip
\noindent\textit{Allowed standard prerequisites (editorial).}\begin{itemize}
\item Quasi-coherence of the pullback of a quasi-coherent module. Same chapter, the lemma that the pullback of a quasi-coherent module along an etale morphism of algebraic spaces is quasi-coherent. Supplies the passage from condition (2) to condition (3).
\item Local presentations of quasi-coherent modules on a site. Modules on Sites, the definition of a globally presented module, the lemma producing local presentations, and the criterion for a module on a site to be quasi-coherent. Supplies the construction of global presentations on the members of an etale covering and the resulting quasi-coherence.
\end{itemize}\medskip
\ifdefined\hypertarget\hypertarget{sourcepointer-1}{}\fi
\noindent\textit{Source reference.} Original equation 03LW: sources/stacks-spaces-properties/pinned-original-source.tex, source lines 4629--4629. Prerequisite pointer invoked at proof line(s) 5636; it names a standard fact of the same chapter that the counted argument uses, and every step of the counted proof built on it is retained verbatim in the delivered body.\par\medskip
\ifdefined\hypertarget\hypertarget{sourcepointer-2}{}\fi
\noindent\textit{Source reference.} Original reference 01BG: sources/stacks-spaces-properties/pinned-original-source.tex, source lines 5461--5461. Prerequisite pointer invoked at proof line(s) 5631; it names a standard fact of the same chapter that the counted argument uses, and every step of the counted proof built on it is retained verbatim in the delivered body.\par\medskip
\ifdefined\hypertarget\hypertarget{sourcepointer-3}{}\fi
\noindent\textit{Source reference.} Original definition 03DE: sources/stacks-chow/tags.txt, source lines 4417--4417. Prerequisite pointer invoked at proof line(s) 5640; it names a standard fact of a sibling chapter that the counted argument uses, and every step of the counted proof built on it is retained verbatim in the delivered body.\par\medskip
\ifdefined\hypertarget\hypertarget{sourcepointer-4}{}\fi
\noindent\textit{Source reference.} Original definition 03DL: sources/stacks-chow/tags.txt, source lines 4424--4424. Prerequisite pointer invoked at proof line(s) 5649; it names a standard fact of a sibling chapter that the counted argument uses, and every step of the counted proof built on it is retained verbatim in the delivered body.\par\medskip
\ifdefined\hypertarget\hypertarget{sourcepointer-5}{}\fi
\noindent\textit{Source reference.} Original lemma 03DN: sources/stacks-chow/tags.txt, source lines 4426--4426. Prerequisite pointer invoked at proof line(s) 5641; it names a standard fact of a sibling chapter that the counted argument uses, and every step of the counted proof built on it is retained verbatim in the delivered body.\par\medskip
\begin{lemma}
\label{lemma-characterize-quasi-coherent}
Let $S$ be a scheme. Let $X$ be an algebraic space over $S$.
Let $\mathcal{F}$ be a sheaf of $\mathcal{O}_X$-modules.
The following are equivalent
\begin{enumerate}
\item $\mathcal{F}$ is a quasi-coherent $\mathcal{O}_X$-module,
\item there exists an \'etale morphism $f : Y \to X$ of
algebraic spaces over $S$ with $|f| : |Y| \to |X|$ surjective
such that $f^*\mathcal{F}$ is quasi-coherent on $Y$,
\item there exists a scheme $U$ and a surjective \'etale morphism
$\varphi : U \to X$ such that $\varphi^*\mathcal{F}$ is a quasi-coherent
$\mathcal{O}_U$-module, and
\item for every affine scheme $U$ and \'etale morphism $\varphi : U \to X$ the
restriction $\varphi^*\mathcal{F}$ is a quasi-coherent $\mathcal{O}_U$-module.
\end{enumerate}
\end{lemma}

\begin{proof}
It is clear that (1) implies (2) by considering $\text{id}_X$.
Assume $f : Y \to X$ is as in (2), and let $V \to Y$ be a surjective
\'etale morphism from a scheme towards $Y$. Then the composition $V \to X$ is
surjective \'etale as well
and by Lemma \ref{lemma-pullback-quasi-coherent} the pullback of $\mathcal{F}$
to $V$ is quasi-coherent as well. Hence we see that (2) implies (3).

\medskip\noindent
Let $U \to X$ be as in (3). Let us use the abuse of notation introduced
in Equation (\ref{equation-restrict-modules}).
As $\mathcal{F}|_{U_\etale}$ is quasi-coherent there exists an
\'etale covering $\{U_i \to U\}$ such that
$\mathcal{F}|_{U_{i, \etale}}$ has a global presentation, see
Modules on Sites, Definition \ref{sites-modules-definition-global} and
Lemma \ref{sites-modules-lemma-local-final-object}.
Let $V \to X$ be an object of $X_\etale$. Since $U \to X$ is
surjective and \'etale, the family of maps $\{U_i \times_X V \to V\}$ is an
\'etale covering
of $V$. Via the morphisms $U_i \times_X V \to U_i$ we can restrict the
global presentations of $\mathcal{F}|_{U_{i, \etale}}$ to get a global
presentation of $\mathcal{F}|_{(U_i \times_X V)_\etale}$
Hence the sheaf $\mathcal{F}$ on $X_\etale$ satisfies the condition of
Modules on Sites, Definition \ref{sites-modules-definition-site-local}
and hence is quasi-coherent.

\medskip\noindent
The equivalence of (3) and (4) comes from the fact that any scheme has
an affine open covering.
\end{proof}

\end{document}
